%! Simplifying Rational Expressions — Unit 1, Lesson 1.7 — Exit Ticket (Answer Key)
\documentclass[10pt]{article}
\usepackage{atda-article}
\usepackage{atda-key}   % pulls in -boxes; do NOT also load -boxes
\newcommand{\TallMath}[1]{$\displaystyle #1\rule[-1.4em]{0pt}{3.2em}$}
\setlength{\workrowsep}{8pt}

\begin{document}

\pageheader{Unit 1, Lesson 1.7}{Exit Ticket --- Answer Key}
% NAMESTRIP: no \namedateperiod here — mirrors the blank. Teacher-only prose goes
% in the lesson plan as \begin{teachernote}[Exit Ticket], never in a key.

\vspace{0.15in}

\begin{notesbox}{Rational Expressions --- On Your Own}
\small
Notes closed. Three items. Factor first, cancel only \emph{factors}, and state every excluded value
from the \emph{original} denominator.

\begin{enumerate}[leftmargin=1.6em, itemsep=8pt, topsep=6pt]

\item Simplify and state every excluded value: \ \TallMath{\frac{x^2-25}{x^2+8x+15}}
\begin{work}
  \frac{x^2-25}{x^2+8x+15} &= \frac{(x+5)(x-5)}{(x+5)(x+3)} \\
      &= \frac{x-5}{x+3} \qquad (x \ne -5, \; x \ne -3)
\end{work}

\item Multiply and simplify: \ \TallMath{\frac{x^2-9}{x+2} \cdot \frac{x^2-4}{x-3}}. State every
excluded value.
\begin{work}
  \frac{(x+3)(x-3)}{x+2} \cdot \frac{(x+2)(x-2)}{x-3}
      &= (x+3)(x-2) \qquad (x \ne -2, \; x \ne 3)
\end{work}

\item A student simplifies \ \TallMath{\frac{x+6}{x+2}} \ to $\dfrac62 = 3$. Test the claim at
$x=2$, say in one sentence what the student did wrong, and give the correct simplest form.
\begin{work}
  \text{at } x=2: \; \frac{2+6}{2+2} &= \frac{8}{4} = 2, \quad \text{not } 3
\end{work}
\ansline{The $x$'s are terms, not factors --- nothing cancels, so $\frac{x+6}{x+2}$ is already simplest.}

\end{enumerate}
\end{notesbox}

\end{document}
